%!TEX root = ../main.tex

\section{Introduction}

%\subsection{Rationale}

%The idea of context-aware computing was first introduced by Schilit et al. in 1994, defining it as a type of software that can dynamically adjust itself to its spatial and temporal environment, as well as the surrounding individuals and objects\cite{schilit1994context}. Dey \cite{dey2001understanding}introduced an alternative task-oriented definition, describing context awareness as using information and offering services relevant to the particular tasks of a person. Today, context awareness is a field of study and a cornerstone of ubiquitous computing.

%There are many approaches to the classification of contextual data used in the literature, depending mainly on the purposes of analysis or theoretical approaches involved in the investigation of the issue \cite{mostefaoui2004contextaware}. In their seminal article on the concept of context-aware computing, Schilit et al. \cite{schilit1994context} defined the three basic types of context – user, computing, and physical contexts represented respectively by a user's profile, geographical coordinates, and computing environment, and ambient lighting and noise level.

%Focusing on context parameters that are especially important for mobile applications, some early research \cite{schilit2000customizingmobile} suggested considering four major categories of context – location, identity, time, and activity as the most basic metrics that make it possible to conclude more specific contextual factors. Thus, identifying the identity of a person allows inferring his preferences, contacts, and relations with others, and the identification of his location allows analyzing potential activities in a place and the presence of other persons or objects around him.

%Context may also be classified on the basis of user's subjective perceptions and sensations. In this case, Gwizdka \cite{gwizdka2000context} distinguished internal and external context. The former implies the identification of mental and emotional states of a person, and includes such factors as task state, personal events, and communication logs, while the latter relates to the conditions of the outer world and involves the location of a person, time, and proximity to other people and devices.

%In terms of HCI theory within location-based systems, Petrelli et al. \cite{petrelli2000modeling} distinguished material and social context. Material context refers to the physical space and existing technological resources (type of a device used, network availability, or GPS), and social context relates to personal characteristics of a user, his interest, and the current context, such as whether he is alone or with somebody else.

%***** ***** ***** ***** ***** *****

The rapid growth of ubiquitous computing and consumer-grade smart devices has fundamentally transformed the landscape of digital mental health systems, as continuous data streams from smartphones and wearables enable real-time assessment and behavioral inference, while the anywhere-anytime paradigm facilitates timely, personalized interventions, giving rise to novel mental care strategies, such as Ecological Momentary Assessments (EMA)\cite{fortuna2022peersupportedema}, Ecological Momentary Interventions (EMI)\cite{balaskas2021ecological}, and more recently, Digital Phenotyping \cite{bufano2023digital} and Just-in-Time Adaptive Interventions (JITAI)\cite{vangenugten2025jitai}. These methodologies complement conventional approaches, and generally aim to ubiquitously capture user-relevant data, analyze and process them into actionable insights and provide real-time, highly granular psychological support beyond traditional clinical settings. 

In all these strategies, context plays a central, albeit poorly specified, role. Indeed, the very notion of context has been widely discussed by the research community. Dey~\cite{dey2001understanding} defines context as any information that characterizes the situation of an entity, where the entity can be a person, a place, or an object relevant to a user, a system, and their interaction. Coutaz and colleagues~\cite{coutaz2005} recognized the challenges in modeling context in a structured yet flexible manner. The authors distinguish between \textit{context-as-process} and \textit{context-as-state}, the latter reflecting the need for context-aware applications to define and store contextual information in a machine-understandable way. This dichotomy is echoed a decade later by Alegre and colleagues~\cite{alegre2016}, who note that context-aware applications can be approached from a phenomenological or a positivist perspective. The former, closer to psychology, explains context in complex human behavior and activity, as an \textit{interaction problem} (i.e., \textit{context-as-process}), in which context varies depending on the situation, the person, and the surrounding social and environmental circumstances. Empirically, this was supported by an exploration study on context-aware applications for mental health intervention~\cite{miralles2019}, where participating psychologists cast doubt on which features of context were most relevant to each type of intervention, as this strongly depended on the person and the mental health disorder at hand. The latter, an empirical tradition of positivism that is closer to computer science, traditionally treat context as a \textit{representation problem}, i.e., \textit{context-as-state}. 

Therefore, we acknowledge the ambiguity of context and, throughout this work, we adhere, as a starting point, to the categorization of Han et al.~\cite{han2008contextaware}, who provided a theoretical framework dividing context into physical (environmental real-world features), internal (internal user-related features), and social (human interaction-based) aspects. In this review, we will further enrich these initial categories through an inductive analysis of the eligible studies. 

Within this categorization, active data acquisition approaches -- particularly notification-delivered questionnaires (i.e., self-reporting) or in-app assessments (e.g., cognitive tasks, structured or freestyle note taking or journaling) -- have traditionally been, and continue to be, predominant \cite{kraft2024mobile, oudin2023digital}. They typically emphasize capturing the patient's internal context~\cite{stone2023evaluation}: emotional~\cite{hickey2021wearablemental} (e.g., mood, stress level) and cognitive~\cite{varandas2022cognitivefatigue} (e.g., attention, mental fatigue) states, (other) mental health-related symptoms~\cite{sheikh2021passivesensing} (e.g., behavioral or physical), psychological functioning~\cite{binmorshed2021selfesteem} (e.g., subjective well-being, self-perception, motivation) and lifestyle factors~\cite{sanudo2020covidactivity} (e.g., sleep, physical activity, daily activities). To complement internal context, these approaches are also used to capture external context (e.g., social context, activities, environment), yet their reliance on repeated user input may be burdensome and limit the breadth and frequency of captured contextual information~\cite{langener2023capturing}. 

In consequence, and driven by the proliferation of modern mobile devices, active sensing has been gradually complemented with passive sensing approaches (i.e., background sensing), whereby the built-in sensors of mobile devices were initially exploited~\cite{miralles2020smartphone, kraft2024mobile}, followed by sensors embedded in wearables devices (e.g., sport bands, smartwatches, smart rings)~\cite{hassan2025utility, jung2025key}. Passive sensing reduces user effort and reporting bias, and provides objective, continuous and unobtrusive monitoring~\cite{shen2025passive}. As such, beyond supplementing and enriching information about the user's internal context (e.g., tracking internal physiological and biological metrics, such as heart rate variability (HRV), electrodermal activity (EDA) and sleep architecture~\cite{sheikh2021passivesensing}% to detect and/or infer physiological stimuli or symptoms\cite{lanata2015complexity}
), passive sensing offers additional continuous and objective insights into the external context, which shape an individual's experiences and behaviors, and may provide insights in the situational causes of mental distress~\cite{onnela2016digitalphenotyping}. % This includes physical dimensions (e.g., physical and semantic location, and noise and ambient light exposure), social dimensions (e.g., social interactions, proximity networks, communication frequencies), and computing/temporal dimensions (e.g., device usage telemetry, temporal behavioral rhythms). 
This notion of digital phenotyping has itself evolved: although the term originally referred almost exclusively to the use of smartphones for data capture~\cite{onnela2016digitalphenotyping}, the concept has progressively expanded~\cite{deboer2023} to account for a wide range of mobile devices, as wearable technology has advanced considerably in functional diversity~\cite{ometov2021}. Today, digital phenotyping seeks to quantify the individual human phenotype (behaviour, physiology) by collecting, processing and analyzing multi-device and multi-modal active and passive data -- encompassing both internal and external contextual factors -- hereby transforming raw digital signals into actionable digital (bio)markers that serve as indicators of symptoms and states, particularly in the mental health domain~\cite{bufano2023digital, moura2023digital}.
%Central to the subsequent delivery of personalized, just-in-time adaptive interventions is the integration of internal and external contextual factors to inform and drive therapeutic responses. 

This synthesis of internal and external signals is, in turn, central to the effective delivery of psychological interventions: integrating an individual's internal context and external context is what allows therapeutic responses to be informed, driven and dynamically tailored to evolving psychological needs and ecological circumstances. Beyond improving the identification and quantification of potential causal triggers, external contextual information enables interventions to be dynamically adapted based on the individual's ecological context, such as location (e.g., absolute, relative location, proximity), temporal factors (e.g., time of day, weekday/weekend, holiday), activity/behavioral context (i.e., engaged activities, daily routines), social circumstances (e.g., being alone, social interactions), situational (e.g., at home, at work, in a public place) or environmental (e.g., weather, noise levels) conditions~\cite{bufano2023digital, onnela2016digitalphenotyping}. 

Despite this potential and emerging integrative research, there remains a divide between digital phenotyping, which primarily focuses on uncovering digital markers for monitoring and prediction~\cite{choi2024digital}, and the integration of passive sensing and resulting markers into personalized, adaptive mental health interventions~\cite{bosch2025integration, choi2024digital}. State-of-the-art JITAI systems still rely primarily on self-reported data, yet with a growing body of literature reporting passive sensing data~\cite{hsu2025personalized, somkuwar2026transforming}, even though advanced concepts, such as (derived) individuals’ states and external contexts, are hardly used~\cite{van2025beyond}. 

% Further opportunities exist to harness social relationships (e.g., multi-person interventions), context-aware multi-modal deliveries \cite{nahum2025just}.

%However, the paradigm in this domain has predominantly focused on passive sensing of internal context\cite{nelson2018smarthome}, specifically internal physiological metrics such as heart rate variability (HRV), electrodermal activity (EDA), and sleep architecture\cite{sheikh2021passivesensing}.

%While tracking internal physiological and biological metrics is highly effective for detecting the presence of physiological stimulus or the manifestation of psychological symptoms\cite{lanata2015complexity}, it suffers from a major critical diagnostic limitation: it identifies that distress is occurring, without understanding the context in which it is occurring. Relying exclusively on internal symptom tracking creates a fundamental blind spot regarding the situational and environmental triggers of mental health disorders. Consequently, there is a distinct "why" gap in current frameworks for ubiquitous computing for mental health.

%To bridge this gap, intervention architectures must shift toward the capture and analysis of external context. External context sensing focuses on the situational causes of distress by continuously modeling the user's environment. This includes physical dimensions ( Physical and semantic location, and noise and ambient light exposure), social dimensions (Social interactions, proximity networks, and communication frequencies), and computing/temporal dimensions (such as device usage telemetry and temporal behavioral rhythms). The integration of multidimensional context variables enables algorithmic systems to map the events that lead to mental health distresses. Transitioning from passive symptom sensing to active, situation-aware monitoring gives a solid foundation for developing symptom-reactive treatments and interventions\cite{teles2017situationawareness, wang2024contextaware}.

% Although this shift is theoretically necessary and beneficial, the extent to which external context variables are effectively extracted, integrated, and utilized within current consumer-grade digital health interventions remains fragmented. A structured evaluation of the literature is required to determine how context-aware computing maps external environmental factors to mental health outcomes.

%Several reviews have examined the field of digital phenotyping and mobile health technology in the past few years. This systematic review, however, occupies a unique methodological niche by shifting the focus from passive monitoring to context-aware intervention. Current literature continues to focus primarily on downstream diagnostic prediction, longitudinal symptom tracking, or limited device ecosystems. For example, Lee et al. (2023) concentrate on digital phenotyping to detect and predict general patterns of behavior, while Busshart et al. (2026) concentrate on identifying the passive sensor streams and active self-reports that best follow the severity of depressive symptoms over time. Likewise, Choi et al. (2024) also evaluate digital phenotyping for stress, anxiety and mild depression, but limit their review strictly to inbuilt smartphone sensors in non-clinical adult groups. Moreover, even when the clinical scope of the reviews is broad, the focus is often overwhelmingly on observation rather than treatment. For example, Dumas et al. (2026) include both psychiatric and non-psychiatric physical or neurological conditions (e.g., Parkinson’s disease, sleep apnea), but the intervention focus is absent, with intervention delivery present only in a minuscule fraction (3\%, 2 out of 65 studies). We make a distinctive contribution by limiting our scope to digital mental health interventions that provide active real-time support, as opposed to these paradigms that are high on tracking and exclusive to phenotyping. Moreover, our eligibility criteria specifically restrict the target population to those with recognized mental health disorders based on DSM-5 guidelines, unlike previous studies that mostly used non-clinical samples or pooled unrelated physical diseases. Instead, we systematically review and classify how a broad hardware ecosystem, including smartphones and consumer-grade wearables, is used to capture a rich, multimodal representation of physical, social, and computing/temporal environments to inform the triggering and tailoring of Just-in-Time Adaptive Interventions (JITAIs) and Ecological Momentary Interventions (EMIs), instead of restricting the use of internal physiological signals or single smartphone sensor configurations.

%\subsection{Objectives}
This systematic review aims to categorize and analyze the use of external contextual factors in digital mental health interventions for recognized mental disorders, delivered through consumer smart devices and wearable technologies. Hereby, this review seeks to systematically map how and which external context information is being captured, and how it is exploited to define context- and situation-aware intervention models. For the obtained body of research, we furthermore determine maturity and degree of evaluation, and finally identify current gaps and future research opportunities. %We furthermore aim to identify technological capabilities and gaps. %\cite{adler2024actionablesensing}.
Specifically, this review addresses the following research questions:
\begin{enumerate}[label=Q\arabic*, itemsep=0pt]   
    \item Which mental health disorders are targeted by context-aware interventions?
    \item What hardware platforms were used to implement these interventions?
    \item What types of external contextual data are captured for these interventions?
    \item What sensing technologies and methods are used to collect external context data?
    \item How is the collected external context data exploited and which specific mechanisms are used, in intervention delivery?
    \item What evaluation designs and dimensions are used to assess the interventions?
\end{enumerate}

%\begin{figure}[H]
%    \centering
%    \includegraphics[width=0.95\textwidth]%{sections/charts/Diagram.png}
%\end{figure}